\section{Cylinder geometry, complete spectrum, and the Green weight}
\label{sec:geometry}

\subsection{The geometric square}

Let $\K=L^2(S^2,\Sigma;\dd\Omega)$ be the unit-sphere spinor space.
Use the self-adjoint convention $D=-\ii\not{\nabla}$ and let
$\Gamma$ be parallel chirality. Then
\begin{equation}
 \Gamma^*=\Gamma,\qquad \Gamma^2=I,\qquad
 \Gamma D=-D\Gamma,
 \qquad D\chi_{n,a}^{\pm}=\pm n\chi_{n,a}^{\pm},
 \quad 1\le a\le2n,\quad n\ge1.
 \label{eq:sphere-spectrum}
\end{equation}
We take $\chi^-_{n,a}=\Gamma\chi^+_{n,a}$ and write $A=D^2$.
The projection $P_n$ onto the complete $A=n^2$ shell has rank $4n$.
These eigenvalues and multiplicities are the $S^2$ specialization of the
spinor spectrum in \cite[eqs.~(3.40), (3.54)]{CamporesiHiguchi1996}.

Consider the complete cylinder
\begin{equation}
 C=\R\times S^2_{1/2},\qquad
 g_C=\dd v^2+\tfrac14g_{S^2},\qquad
 \dd\mathrm{vol}_C=\tfrac14\dd v\dd\Omega.
\end{equation}
In the product-measure representation its geometric Dirac operator is
\begin{equation}
 \Dc=-\ii\Gamma\partial_v+2D,
 \qquad L_0=\Dc^2=-\partial_v^2+4A.
 \label{eq:cylinder-square}
\end{equation}
The cross terms cancel by the anticommutation in
\eqref{eq:sphere-spectrum}. Multiplication by $2$ is the unitary map from
product-measure spinors to metric-volume spinors. It does not change the
differential expression.

The cylinder scalar curvature is $8$. The Lichnerowicz identity therefore
also reads $\Dc^2=\nabla^*\nabla+2$. This is the same operator as
\eqref{eq:cylinder-square}: the curvature contribution is already inside
the Dirac square, and adding a further $2$ would change the model.
The slash and curvature conventions agree with
\cite[eqs.~(5.35)--(5.36)]{CamporesiHiguchi1996} after multiplication by
$-\ii$ to make the Dirac operator self-adjoint.

For a length scale $\ell_0>0$ and mass $m_{\mathrm{phys}}>0$, the supplied
energy convention is
\begin{equation}
 H_{\mathrm{kin}}=\varepsilon L_0,
 \qquad \varepsilon=\frac{\hbar^2}{2m_{\mathrm{phys}}\ell_0^2},
 \qquad \tau=\frac{\varepsilon t_E}{\hbar}.
 \label{eq:physical-scale}
\end{equation}
Thus $\ee^{-\tau L_0}$ is Euclidean evolution. Real laboratory time would
instead give $\ee^{-\ii\varepsilon t_{\mathrm{lab}}L_0/\hbar}$.

\subsection{Self-adjoint realization without a radial cutoff}

Use the unitary full-line Fourier convention
\begin{equation}
 \widehat f(k)=(2\pi)^{-1/2}\int_\R\ee^{-\ii kv}f(v)\dd v.
 \label{eq:fourier}
\end{equation}
In each pair $(\chi^+_{n,a},\chi^-_{n,a})$ the first-order multiplier is
\begin{equation}
 M_n(k)=\begin{pmatrix}2n&k\\k&-2n\end{pmatrix},\qquad
 E_n(k)=\sqrt{k^2+4n^2},\qquad
 \Pi_\pm(n,k)=\tfrac12\left(I\pm\frac{M_n(k)}{E_n(k)}\right).
 \label{eq:signed-block}
\end{equation}
Since $M_n(k)^2=\omega_n(k)I$ with $\omega_n(k)=k^2+4n^2$,
the maximal domains on $L^2(\R;\K)$ are
\begin{align}
 \Dom\Dc=\Dom L_0^{1/2}
 &=\left\{f:\sum_n\int_\R\omega_n(k)
           \norm{P_n\widehat f(k)}^2\dd k<\infty\right\},
 \label{eq:dirac-domain}\\
 \Dom L_0^p
 &=\left\{f:\sum_n\int_\R\omega_n(k)^{2p}
           \norm{P_n\widehat f(k)}^2\dd k<\infty\right\},
 \qquad p\ge0.
 \label{eq:square-domain}
\end{align}

\begin{proposition}[Complete realization]
\label{prop:complete-realization}
The multipliers above define self-adjoint operators. Finite complete-shell
sums with radial coefficients in $C_c^\infty(\R)$ are cores for $\Dc$
and every fixed finite integer power of $L_0$. Their spectra are
\begin{equation}
 \sigma(\Dc)=(-\infty,-2]\cup[2,\infty),\qquad
 \sigma(L_0)=[4,\infty),
 \label{eq:spectra}
\end{equation}
and are purely absolutely continuous in the radial spectral variable.
\end{proposition}

\begin{proof}
Hermitian matrix multiplication is symmetric on its maximal domain.
Testing a vector in the adjoint domain against compactly supported
spectral functions in finitely many channels forces its image under
$M_n(k)$ to be square-integrable. This gives exactly
\eqref{eq:dirac-domain}; the scalar argument gives
\eqref{eq:square-domain}. Hence the maximal operators equal their
adjoints.

Projection onto $\sum_{n\le N}P_n$ converges in each displayed graph norm
by dominated convergence and commutes with the signed blocks. On this
finite angular space the first-order graph norm is equivalent to the
radial $H^1$ norm, and the graph norm of $L_0^p$ to $H^{2p}$ for integer
$p$. Smooth cutoffs and mollification give the stated cores. No radial
boundary condition is introduced.

The eigenvalues of \eqref{eq:signed-block} are $\pm E_n(k)$, each of
physical multiplicity $2n$; the square has multiplicity $4n$. Their
essential ranges give \eqref{eq:spectra}. On each open half-line in $k$,
the dispersions are smooth and strictly monotone with nonzero derivative.
Pushing an integrable spectral density through them gives an absolutely
continuous measure. The exceptional threshold preimages are singletons
and have measure zero. Countable sums preserve absolute continuity.
\end{proof}

In particular, a fixed radial plane wave, cosine, or threshold value is
not an $L^2$ eigenstate. Restricting to even radial functions reduces the
square, but not the first-order operator: its radial derivative changes
parity. Complete shells are the natural cutoffs when both signs of the
Dirac operator are retained.

\subsection{The Green diagonal and multiplicity cancellation}

On shell $n$, the inverse of $-\partial_v^2+4n^2$ has kernel
\begin{equation}
 g_n(v,w)=\frac{\ee^{-2n|v-w|}}{4n}.
 \label{eq:green-kernel}
\end{equation}
Indeed, away from $v=w$ it solves the homogeneous equation, while its
first derivative has jump $-1$, producing the required unit delta
under $-\partial_v^2$. Equivalently,
\begin{equation}
 g_n(v,v)=\int_\R\frac{\dd k}{2\pi(k^2+4n^2)}=\frac1{4n}.
 \label{eq:green-diagonal}
\end{equation}
This is the complete radial integral. Evaluation of the resolvent at
$k=0$ would give $1/(4n^2)$ and would not yield this weight.

The angular operator supplied by the diagonal is
\begin{equation}
 W=(4\sqrt A)^{-1}=\sum_{n\ge1}\frac1{4n}P_n.
 \label{eq:green-weight}
\end{equation}
It is bounded and positive but not trace class. For every $\beta>0$,
however,
\begin{equation}
 \Tr_\K(W\ee^{-\beta A})
 =\sum_{n\ge1}(4n)\frac1{4n}\ee^{-\beta n^2}
 =Q_1(\beta).
 \label{eq:theta-trace}
\end{equation}
This exact cancellation is the link between the cylinder geometry and
the unit coefficients of the theta series. It retains all shells and
both radial directions. It is not the trace of the full cylinder heat
operator, which has infinite radial volume.
